\section{2026 年 11 月核验增补}

本节按论文首次提交时间排列。每项方法附有可核验的外部入选依据。

% METHOD: 2609.35612
\paragraph{DIBR-GS：用逐步合成的深度伪视图补足遥感观测}
2026.09~\cite{2609_35612}。DIBR-GS 先把单目估计的深度与 SfM 稀疏点对齐，并融合多视图来初始化锚点。表示端加入来自参考图像的跨视角外观特征，训练端逐步把观测深度重投影到未观测视点，为高斯提供附加一致性约束。新的伪视图传递已有几何证据，并未增加真实采集。

作者在 LEVIR-NVS 上比较遥感新视图合成，观察到结构边界和细节的改进。框架仍需逐场景优化，伪视图质量还取决于单目深度和稀疏几何的对齐，不能直接转用为前馈重建。

\textbf{入选依据} \href{https://github.com/kanehub/dibr-gs}{公开实现}已找到与论文对应的代码文件。
